\documentclass[11pt]{article}
\usepackage[margin=1in]{geometry}
\usepackage{amsmath,amssymb,amsthm}
\usepackage{xurl}
\usepackage{hyperref}
\sloppy
\let\oldus\_ \renewcommand{\_}{\oldus\allowbreak}
\newtheorem{theorem}{Theorem}
\newtheorem{lemma}[theorem]{Lemma}
\theoremstyle{definition}
\newtheorem{definition}[theorem]{Definition}
\newcommand{\K}{\mathbb{K}}
\newcommand{\ip}[2]{\langle #1,#2\rangle}

\title{Proof of conjecture 00000000960:\\ $\ell^2(\mathbb{R})$ has no Schauder basis, but every separable subspace has one}
\author{Jackmeson1}
\date{2026-10-05}

\begin{document}
\maketitle

\section{The conjecture and its reading}

\textbf{Statement (English, verbatim).} Conjecture: There exists a Banach space with no Schauder
basis, each of whose separable subspaces has one (the Enflo direction of the basis problem).
The Chinese text says the same: there exists a Banach space with no Schauder basis, every
separable subspace of which has a basis.

\textbf{Reading.} The claim is the existential statement
\[
  \exists\ \text{Banach space } X:\quad X \text{ has no Schauder basis} \ \wedge\
  \forall\, V\le X \text{ separable}:\ V \text{ has a Schauder basis}.
\]
If such an $X$ were separable, it would be a separable subspace of itself and would have a basis
and also not have one. So any witness is non-separable. We prove the statement with
$X=\ell^2(\mathbb{R},\K)$, $\K=\mathbb{R}$ or $\mathbb{C}$.

\textbf{Conventions (all used in the proof and the Lean code).}
\begin{enumerate}
\item \emph{Schauder basis} means the classical, sequence-indexed notion. A Schauder basis of a
normed space $Y$ is a sequence $(e_n)_{n\in\mathbb{N}}$ in $Y$ with continuous linear functionals
$(f_n)$ such that $f_i(e_j)=\delta_{ij}$ and $\sum_{i<n} f_i(x)e_i\to x$ as $n\to\infty$ for every
$x\in Y$. This is Mathlib's \texttt{SchauderBasis}~$\K$~$Y$ (an abbreviation of
\texttt{GeneralSchauderBasis} $\mathbb{N}$ $\K$ $Y$ with the conditional summation filter),
\nolinkurl{Mathlib/Analysis/Normed/Module/Bases.lean}. Some authors also call an \emph{uncountable}
family a Schauder basis (Source~1 says so). Under that convention $\ell^2(\mathbb{R})$ has one, its
orthonormal basis, and this proof does not apply. See Section~4.
\item \emph{No basis} is proved in a strong form. $X$ has no Mathlib \texttt{SchauderBasis}. More
generally, $X$ has no \texttt{GeneralSchauderBasis} indexed by any countable type along any
non-trivial summation filter. It also has no sequence $(e_n)$ such that every $x$ is the limit of
the partial sums of some series $\sum a_n e_n$. That last form needs neither uniqueness nor
continuity of the coefficients, so it also rules out the definition in Source~1 (a sequence with
unique expansions), with or without zero-padded finite bases.
\item \emph{Separable subspace} means any linear subspace $V\le X$ (closed \emph{or not}) whose
underlying set is separable, with the inherited norm. We do not restrict to closed subspaces.
\item \emph{Finite-dimensional subspaces.} An $\mathbb{N}$-indexed Schauder basis is linearly
independent, so a finite-dimensional space never has one. For finite-dimensional $V$ we give a
finite basis: a \texttt{GeneralSchauderBasis} indexed by \texttt{Fin}~$(\dim V)$ (biorthogonal,
with $x=\sum_{i<\dim V} f_i(x)e_i$). For infinite-dimensional $V$ we give a classical
$\mathbb{N}$-indexed Schauder basis.
\end{enumerate}

\textbf{The parenthetical.} ``The Enflo direction of the basis problem'' refers to Banach's basis
problem. Source~1: ``A Banach space with a Schauder basis is necessarily separable, but the converse
is false. The basis problem is the question asked by Banach, whether every separable Banach space has
a Schauder basis. This was negatively answered by Per Enflo who constructed a reflexive and
separable Banach space failing the approximation property, thus a space without a Schauder
basis.'' (Enflo, Acta Math.\ 130 (1973) 309--317, \url{https://doi.org/10.1007/bf02392270},
Source~2.) Enflo's space is separable, so it is
a separable subspace of itself without a basis. It is therefore \emph{not} a witness for the
stated existential, which is a different and easier statement. We prove the statement as written.

\section{Proof}

Let $X=\ell^2(\mathbb{R},\K)$, the Hilbert space of square-summable families
$(x_t)_{t\in\mathbb{R}}$. It is complete.

\begin{lemma}[countable bases force separability]\label{l:sep}
Let $Y$ be a normed space over $\K$ and $s\subseteq Y$ a countable set. If every $x\in Y$ lies in
the closure of $\operatorname{span}s$, then $Y$ is separable. Hence a space with a Schauder basis,
with a countably indexed generalized basis, or with a sequence $(e_n)$ such that every $x$ is a
limit of partial sums $\sum_{i<n}a_ie_i$, is separable.
\end{lemma}
\begin{proof}
The span of a countable set over the separable field $\K$ is separable, and so is its closure,
which is $Y$. In each of the three cases every $x$ is a limit of finite linear combinations of the
countably many basis vectors.
\end{proof}

\begin{lemma}\label{l:nonsep}
$X$ is not separable.
\end{lemma}
\begin{proof}
Let $(b_t)_{t\in\mathbb{R}}$ be an orthonormal family indexed by $\mathbb{R}$ (Mathlib's default
Hilbert basis of $\ell^2(\mathbb{R},\K)$). For $s\ne t$,
$\|b_s-b_t\|^2=\|b_s\|^2-2\operatorname{Re}\ip{b_s}{b_t}+\|b_t\|^2=2$, so the open balls
$B(b_t,1/2)$ are pairwise disjoint. In a separable space a family of pairwise disjoint non-empty open
sets is countable. But $\mathbb{R}$ is uncountable.
\end{proof}

By Lemmas~\ref{l:sep} and~\ref{l:nonsep}, $X$ has no Schauder basis in any of the senses of
Convention~2.

\begin{lemma}[expansion in an incomplete space]\label{l:exp}
Let $E$ be an inner product space, not necessarily complete, and let $(v_i)_{i\in I}$ be orthonormal
with dense span. Then $x=\sum_i\ip{v_i}{x}v_i$, with unconditional convergence, for every $x\in E$.
\end{lemma}
\begin{proof}
Let $j:E\to\widehat E$ be the isometric embedding into the completion. The family $(j v_i)$ is
orthonormal, and its span $j(\operatorname{span}\{v_i\})$ is dense, because $j$ is continuous with
dense range. So it is a Hilbert basis of $\widehat E$, and $jx=\sum_i\ip{jv_i}{jx}jv_i
=\sum_i\ip{v_i}{x}jv_i$. Since $j$ is a linear isometric embedding, the finite partial sums
converge to $x$ in $E$.
\end{proof}

\begin{lemma}\label{l:sepbasis}
Let $E$ be a separable inner product space, not necessarily complete. If $E$ is
infinite-dimensional, it has a classical Schauder basis. If it is finite-dimensional, it has a
finite basis of length $\dim E$.
\end{lemma}
\begin{proof}
Finite-dimensional case: take an orthonormal basis $(v_i)_{i<\dim E}$ and the coordinates
$f_i=\ip{v_i}{\cdot}$. Infinite-dimensional case: take a dense sequence $(x_n)$ and apply the
normalized Gram--Schmidt process. Discard the indices where the output is $0$. This leaves an
orthonormal family $(v_i)_{i\in S}$, $S\subseteq\mathbb{N}$, with
$\operatorname{span}\{v_i\}=\operatorname{span}\{x_n\}$, which is dense. If $S$ were finite, the
expansion of Lemma~\ref{l:exp} would be a finite sum and $E$ would be finite-dimensional. So $S$
is countably infinite. Enumerate it by $\mathbb{N}$, which gives an orthonormal sequence $(e_n)$
with dense span. Let $f_n=\ip{e_n}{\cdot}$. These are continuous, and $f_m(e_n)=\delta_{mn}$. By
Lemma~\ref{l:exp}, $\sum_n f_n(x)e_n$ converges unconditionally to $x$, so in particular its
partial sums converge to $x$.
\end{proof}

\begin{theorem}
$X=\ell^2(\mathbb{R},\K)$ is a Banach space without a Schauder basis. Every separable linear
subspace $V\le X$, closed or not, has a classical Schauder basis if it is infinite-dimensional and a
finite basis if it is finite-dimensional. Moreover, $X$ has closed, separable, infinite-dimensional
subspaces, so the second clause is not vacuous.
\end{theorem}
\begin{proof}
The first claim is Lemmas~\ref{l:sep} and~\ref{l:nonsep}. For the second, $V$ with the inherited
inner product is a separable inner product space, so Lemma~\ref{l:sepbasis} applies. For the
third, the closure of $\operatorname{span}\{b_n:n\in\mathbb{N}\}$ is closed and separable, and it
contains infinitely many orthonormal, hence linearly independent, vectors.
\end{proof}

\section{Lean formalization and axiom audit}

The file \texttt{lean/Conjecture960/Basic.lean} (256 lines, \texttt{import Mathlib} only) is
byte-identical to the checked file. It uses the namespace \texttt{C960}. Below, \texttt{K}
stands for the Lean scalar variable (a blackboard-bold k in the Lean source). It ranges over
every \texttt{RCLike} field in \texttt{Type}, so $\mathbb{R}$ or $\mathbb{C}$. $X$ is Mathlib's
$\ell^2(\mathbb{R},\K)$, i.e.\ \texttt{lp (fun \_ => K) 2} over the index type $\mathbb{R}$.
\begin{itemize}
\item \texttt{separableSpace\_of\_closure\_span}, \texttt{separableSpace\_of\_series},
\texttt{separableSpace\_of\_generalSchauderBasis}: Lemma~\ref{l:sep}.
\item \texttt{not\_separableSpace\_l2}: Lemma~\ref{l:nonsep}.
\item \texttt{l2\_no\_countable\_basis}: \texttt{IsEmpty (SchauderBasis K X)}, no countably
indexed \texttt{GeneralSchauderBasis}, and no representing sequence.
\item \texttt{hasSum\_of\_orthonormal\_dense}, \texttt{onbBasis}: Lemma~\ref{l:exp} and the
resulting basis $(v_i,\ip{v_i}{\cdot})$. \texttt{toSchauder}: an $\mathbb{N}$-indexed
unconditional basis is a classical one.
\item \texttt{exists\_orthonormal\_dense} (Gram--Schmidt, Mathlib's \texttt{gramSchmidtNormed}),
\texttt{finiteDimensional\_of\_basis}, \texttt{finite\_has\_basis},
\texttt{separable\_has\_schauderBasis}: Lemma~\ref{l:sepbasis}.
\item \texttt{l2\_separable\_subspace\_has\_basis}, \texttt{l2\_exists\_closed\_separable\_infinite}.
\item The main theorem \texttt{conjecture960} states that there exist \texttt{X : Type},
\texttt{NormedAddCommGroup X} and \texttt{NormedSpace K X} such that:
  \begin{itemize}
  \item \texttt{CompleteSpace X} and \texttt{not (SeparableSpace X)};
  \item \texttt{IsEmpty (SchauderBasis K X)};
  \item for every countable \texttt{iota : Type} and every \texttt{SummationFilter} $L$ on it with
  \texttt{L.NeBot}, \texttt{IsEmpty (GeneralSchauderBasis iota K X L)};
  \item there is no $e:\mathbb{N}\to X$ such that every $x$ has scalars $a:\mathbb{N}\to\K$ with
  $\sum_{i<n}a_ie_i\to x$ (\texttt{Tendsto \ldots\ atTop (nhds x)});
  \item for every \texttt{V : Submodule K X} with \texttt{IsSeparable (V : Set X)},
  \texttt{FiniteDimensional K V} implies
  \texttt{Nonempty (UnconditionalSchauderBasis (Fin (Module.finrank K V)) K V)}, and
  \texttt{not (FiniteDimensional K V)} implies \texttt{Nonempty (SchauderBasis K V)};
  \item some \texttt{V : Submodule K X} is closed, separable and not finite-dimensional.
  \end{itemize}
\item \texttt{conjecture960\_real} restates the real case in the form closest to the wording: a
complete real normed space with \texttt{IsEmpty (SchauderBasis R X)} in which every separable,
infinite-dimensional subspace has a \texttt{SchauderBasis}.
\end{itemize}
\textbf{Axiom audit.} \texttt{\#print axioms} for \texttt{C960.conjecture960} and
\texttt{C960.conjecture960\_real} reports exactly \texttt{[propext, Classical.choice,
Quot.sound]}. There is no \texttt{sorry}, no \texttt{native\_decide} and no added axiom.
\texttt{lake build} succeeds (Lean v4.33.1, Mathlib v4.33.1).

\section{What is not covered}
Under the convention that admits uncountable, linearly ordered ``Schauder bases'' (Source~1:
``some authors define Schauder bases to be countable (as above), while others use the term to
include uncountable bases''), $\ell^2(\mathbb{R})$ has a basis, its orthonormal basis, and is not a
witness. We do not address that reading. The conjecture says ``Schauder basis'', and the definition
quoted in Source~1 (``A Schauder basis is a sequence $\{b_n\}$ \ldots'') is sequence-indexed, so we
use the countable notion.

\section*{Sources}
\begin{enumerate}
\item Wikipedia, ``Schauder basis'',
\url{https://en.wikipedia.org/wiki/Schauder_basis} (retrieved 2026-10-05). Quotes used: ``A
Schauder basis is a sequence $\{b_n\}$ of elements of $V$ such that for every element $v\in V$
there exists a unique sequence $\{\alpha_n\}$ of scalars in $F$ so that $v=\sum\alpha_nb_n$'';
``Note that some authors define Schauder bases to be countable (as above), while others use the
term to include uncountable bases''; and the basis-problem passage quoted in Section~1.
\item P.~Enflo, ``A counterexample to the approximation problem in Banach spaces'', Acta
Mathematica 130 (1973) 309--317, \url{https://doi.org/10.1007/bf02392270} (Crossref record
retrieved 2026-10-05).
\item Mathlib: \nolinkurl{Mathlib/Analysis/Normed/Module/Bases.lean} (\texttt{SchauderBasis},
\texttt{GeneralSchauderBasis}), \nolinkurl{Mathlib/Analysis/InnerProductSpace/l2Space.lean}
(\texttt{HilbertBasis.mk}, \texttt{hasSum\_repr}),
\nolinkurl{Mathlib/Analysis/InnerProductSpace/GramSchmidtOrtho.lean}.
\end{enumerate}

\end{document}
